% Supplementary Information -- English version (v3, 2026-08-14). Compile with pdflatex.
\documentclass[11pt,a4paper]{article}
\usepackage[a4paper,margin=2.4cm]{geometry}
\usepackage{amsmath,amssymb,amsthm,bm}
\usepackage{booktabs,array}
\usepackage{xcolor}
\usepackage[hidelinks]{hyperref}
\newtheorem{theorem}{Theorem}
\newtheorem{corollary}{Corollary}
\renewcommand{\proofname}{Proof}
\newcommand{\bu}{\bm{u}}
\newcommand{\cH}{\mathcal{H}}
\newcommand{\Ex}{E}
\newcommand{\inner}[2]{\langle #1, #2 \rangle}

\title{\bfseries Supplementary Information\\[4pt]
\large Architecture-embedded physical constraints enable stable long-horizon neural PDE prediction}
\author{Yuze Chen}
\date{}

\begin{document}
\maketitle

\section*{SI-A \quad Full proofs of the theorems}

\paragraph{Notation.} $\bu_t$ is the state at step $t$; $\Ex(\bu)=\tfrac12\langle\|\bu\|^2\rangle_\Omega$; $\mathbb{P}$ denotes the Leray spectral projection (Eq.~(3) of the main text); the envelope constants $c_d,c_p$ follow the calibration protocol of the main text.

\begin{theorem}[Divergence-free output; Theorem~1 of the main text]
Suppose the hard anchor contains the Leray projection $\mathbb{P}$ and every step's correction exits through the projection. Then for any input $\bu_t$, the rollout output satisfies $\nabla\!\cdot\!\bu_{t+1}=0$ (to spectral accuracy).
\end{theorem}

\begin{proof}
A periodic $L^2$ vector field admits the Helmholtz decomposition $\tilde\bu=\mathbb{P}\tilde\bu+\nabla q$. For any $\bm{k}\neq\bm{0}$,
\[
\bm{k}\cdot\widehat{\mathbb{P}\tilde\bu}(\bm{k})
=\bm{k}\cdot\widehat{\tilde\bu}(\bm{k})-\frac{\bm{k}\,\big(\bm{k}^{\mathsf T}\widehat{\tilde\bu}(\bm{k})\big)}{\|\bm{k}\|^2}=0 .
\]
In Fourier space, $\bm{k}\cdot\widehat{\bu}(\bm{k})=0$ for all $\bm{k}\neq\bm{0}$ is equivalent to $\nabla\!\cdot\!\bu=0$ in physical space; the $\bm{k}=\bm{0}$ mode is a constant field with zero divergence. $\mathbb{P}$ is idempotent ($\mathbb{P}^2=\mathbb{P}$, since $\bm{I}-\bm{k}\bm{k}^{\mathsf T}/\|\bm{k}\|^2$ is an orthogonal projection matrix), so any output that passes through $\mathbb{P}$ after arbitrary processing lies in the range of $\mathbb{P}$. Every rollout step is therefore divergence-free.
\end{proof}

\begin{theorem}[Dissipation compatibility; Theorem~2 of the main text]
With the envelope active: (i) in decaying flow each step's output energy satisfies $\Ex(\bu_{t+1})\in[\Ex(\bu_t)-c_d\varepsilon(\bu_t)\Delta t,\ \Ex(\bu_t)]$; (ii) in driven flow it falls inside the two-sided interval of Eq.~(9) of the main text; (iii) in decaying flow the energy sequence is non-increasing and bounded below, hence convergent.
\end{theorem}

\begin{proof}
(i)(ii) By the rescaling construction: when the prediction exceeds the upper bound $\Ex_{\max}$, take $\alpha=\sqrt{\Ex_{\max}/\Ex(\bu^{\mathrm{pred}})}<1$; the identity $\Ex(\alpha\bu)=\alpha^2\Ex(\bu)$ (scalar multiplication commutes with the quadratic form) gives $\Ex(\alpha\bu^{\mathrm{pred}})=\Ex_{\max}$; the lower bound is analogous; without a violation, $\alpha=1$. Each step's output energy lands \emph{exactly} inside the closed interval---a constructive identity independent of any network property. (iii) In decaying flow $\Ex_{\max}=\Ex(\bu_t)$ makes the sequence non-increasing; $\Ex\ge0$ bounds it below; a monotone bounded sequence converges.
\end{proof}

\begin{theorem}[Bounded correction / soft landing; Theorem~3 of the main text]
The fused output satisfies $\|\bu_{\mathrm{fused}}-\bu_{\mathrm{anchor}}\|\le R_{\mathrm{eff}}$; at initialisation ($\tau>0$, random parameters) the contraction factor obeys $\lambda<1$; as the residual learner's violation production decays during training, $\lambda\to1$.
\end{theorem}

\begin{proof}
Boundedness: write the untruncated fusion $\bar\bu=(1-w)\bu_{\mathrm{anchor}}+w\bu_{\mathrm{soft}}$; the deviation is $\bar\bu-\bu_{\mathrm{anchor}}=w(\bu_{\mathrm{soft}}-\bu_{\mathrm{anchor}})$. If its norm exceeds $R_{\mathrm{eff}}$, the spherical projection rescales it along the ray to radius $R_{\mathrm{eff}}$, so the inequality holds by construction. Soft landing: the contraction factor is $\lambda=\exp(-\Sigma/\tau)$, where $\Sigma$ is the violation magnitude and $\tau>0$; any non-trivial prediction has $\Sigma>0\Rightarrow\lambda<1$; training drives the residual learner toward physical compatibility, $\Sigma\downarrow0\Rightarrow\lambda\uparrow1$.
\end{proof}

\begin{corollary}[Finite-horizon rollout boundedness; Corollary~1 of the main text]
Suppose the residual learner is Lipschitz in its input (constant $L$, structurally maintained by spectral normalisation), the anchor is non-expansive (both the orthogonal projection and the linear dissipative propagator satisfy $\|\cH(\bu)-\cH(\bm v)\|\le\|\bu-\bm v\|$), and the envelope and gate are both active. Then any $K$-step trajectory is (i) stepwise divergence-free; (ii) monotonically energy-feasible in decaying flow; (iii) uniformly bounded in $\|\bu_t\|$.
\end{corollary}

\begin{proof}
(i)(ii) follow from the stepwise application of Theorems~1 and~2. (iii) By induction: suppose $\|\bu_{t-1}\|\le M_{t-1}$. Non-expansiveness of the anchor gives $\|\cH(\bu_{t-1})\|\le\|\cH(\bm 0)\|+M_{t-1}$; the gate truncation bounds the fused output's deviation from the anchor by $R_{\mathrm{eff}}$; the envelope truncates the energy (hence the norm) to the order of $\sqrt{2\Ex_{\max}|\Omega|}$. Composing the three yields a stepwise bounded recursion $M_t\le M_{t-1}+C$, and boundedness of the energy forces $M_t$ to be uniformly bounded, excluding numerical blow-up.
\end{proof}

\textbf{Testability of the Lipschitz assumption}: spectral normalisation on the hard side structurally maintains upper bounds on the spectral norms of all linear layers during training; per-layer spectral norms can be logged, so the assumption is a checkable numerical property.

\section*{SI-B \quad Point-by-point comparison with structure-preserving numerical methods}

\begin{table}[htbp]
\centering\small
\caption{Correspondence between our architectural priors and classical structure-preserving numerical methods.}
\begin{tabular}{llll}
\toprule
Classical construction & Constrained object & Our counterpart & Guarantee \\
\midrule
Pressure projection (fractional step) & Known discrete NS scheme & Leray projection (Eq.~3) & Theorem 1: divergence-free \\
Energy limiter / monotone scheme & Known flux discretisation & Dissipation envelope (Eqs.~8--10) & Theorem 2: dissipation-compatible \\
CFL-type boundedness condition & Time step & Trust-radius $R_{\mathrm{eff}}$ gating (Eqs.~12--13) & Theorem 3: bounded correction \\
Symplectic / variational integrator & Hamiltonian discretisation & (Not applicable: our systems are dissipative) & --- \\
\bottomrule
\end{tabular}
\end{table}

The key difference: classical methods constrain \emph{known discrete operators}, and their stability proofs rely on the algebraic structure of the scheme; we constrain \emph{a neural network with unanalysable degrees of freedom}, so the guarantees must be constructively independent of network properties---the proofs use only projection idempotency and the scalar scaling identity. The two families do not compete: our target setting is rollout where the true solution is unavailable and data-driven extrapolation is unavoidable; the architectural prior carries structure-preserving ideas into a space where no scheme exists.

\section*{SI-C \quad Data archaeology: a complete record of deprecated data and conclusions}

Self-correction in science should leave a trace. During data verification this study identified and deprecated two legacy datasets together with the conclusions drawn on them:

\paragraph{C.1 KS data (System C).} The legacy \texttt{ks\_L32/L44} files were found by equation-residual verification to satisfy a \textbf{damped variant} $u_t=-uu_x\boldsymbol{+}u_{xx}-u_{xxxx}$ (relative residual $1.85$ against standard KS versus $1.7\times10^{-3}$ against the damped variant), with monotonically decaying energy and no chaos; the legacy \texttt{ks\_L22} file contained 600 frames of pure NaN. The direct cause of the legacy line's ``all length-OOD runs crash'' was NaN data and a $\Delta t$ mismatch ($0.05$ vs $0.01$), not a model property. We also found that the direct-formula ETDRK4's $10^{-12}$ safeguard swallowed the near-zero eigenvalue modes at $L{=}44$ ($\lambda\approx+0.001$), causing numerical blow-up; after switching to the Kassam--Trefethen contour-integral ETDRK4, the one-step integration-consistency relative error is $\sim3\times10^{-8}$. All three length datasets were regenerated with standard KS (burn-in $T{=}50$, 2/3 de-aliasing).

\paragraph{C.2 System B data.} All 8 legacy ``2D'' candidate files were measured to be \textbf{3D velocity fields} ($[T,3,64,64,64]$); the legacy protocol took a mid-plane $(u,v)$ slice and passed it off as 2D NS: the slice divergence ratio is $\approx1.0$ (violating 2D incompressibility) and the vorticity-equation residual is $2$--$24$ (violating 2D NS); the main-system file was even the laminar Kolmogorov profile in its spin-up phase (neither turbulent nor stationary). The legacy protocol also contained IC-statistics fabrication (at $K{=}200$ only 1 IC was actually used while n\_ic$=5$ was reported) and evaluation-start-point leakage into the training segment. All numbers in the legacy \texttt{generality\_system\_b.json} are deprecated.

\paragraph{C.3 Crash-criterion calibration.} The legacy internal baseline line reported ``zero crashes'' under a lenient rL2${}>10$ criterion; its one-sided $1.05\times$/step anti-blow-up net compounds over 200 steps into an effective energy ceiling of $1.05^{200}\approx1.7\times10^{4}$---no constraint at all. This work uniformly adopts the strict criteria of the main text (Crash-criteria table).

\paragraph{C.4 Lessons and rules.} These findings motivated three project rules, now written into the data-generation pipeline: (i) before entering the repository, any dataset must pass an equation-residual check (relative residual $<10^{-2}$) and a one-step integration-consistency check ($<10^{-6}$); (ii) stationarity criteria are calibrated to correlated processes (estimated autocorrelation time $\tau_c$ plus quantile gates on a long reference trajectory)---a white-noise criterion would falsely kill turbulence; (iii) deprecated files are not deleted but archived together with the verification scripts to preserve traceability.

\section*{SI-D \quad Supplementary statistics}

Structure functions $S_2(r)$ and scale-resolved skewness/kurtosis of velocity and vorticity increments (7 models $\times$ 3 IC snapshots, commit \texttt{d96d47c}): our backbone's distributions track the truth most closely at all scales; the intermittency artefact of the gate90 arm (vorticity kurtosis $80.8$) appears consistently at all scales. Per-IC raw rollout trajectories (rL2$(k)$, $E(k)$, $Z(k)$, crash step) are released with the result JSONs.

\end{document}
